\section{Uncertainty, confidence, and missing data}
\label{sec:uncertainty}

\subsection{Ranges, not point estimates}

All quantitative outputs --- temperature reduction, energy savings, greenhouse
gases avoided, cost savings, payback --- are reported as ranges. This is not a
presentational preference; it is what the evidence supports.

The adopted cooling envelopes of \Cref{tab:typologies} span factors of two to
six: street tree planting from \qty{0.5}{\degreeCelsius} to
\qty{3.0}{\degreeCelsius} is a sixfold range, and green roof from
\qty{0.1}{\degreeCelsius} to \qty{1.0}{\degreeCelsius} is tenfold. The
climate-subzone estimates in \textcite{keravec2026} carry uncertainty intervals
ranging from about $\pm\,$\qty{0.6}{\degreeCelsius} to about
$\pm\,$\qty{2.0}{\degreeCelsius}. A single figure drawn from within a spread of
that magnitude would imply a precision that does not exist, and would be
quoted, in practice, without its provenance.

The ranges compound. An energy estimate combines a temperature interval with a
sensitivity interval (\Cref{eq:energy-range}), so its width is the product of
the two --- the worked example in \Cref{sec:worked-example} produces an energy
range spanning roughly an order of magnitude. That width is uncomfortable to
present and it is the honest answer. A user who needs a tighter energy figure
needs a building energy model, which is precisely the boundary
\Cref{sec:not} draws.

\subsection{Confidence levels}
\label{sec:confidence-levels}

Two distinct things are called confidence in this methodology, and separating
them prevents a common confusion.

\textbf{Evidence confidence} is a property of a \emph{typology}, recorded in
the library, expressing how well grounded its adopted values are:

\begin{description}[leftmargin=1.5em, style=sameline, font=\normalfont\itshape]
  \item[High] multiple converging sources, including a synthesis.
  \item[Medium] a synthesis estimate with limited typology-specific
  corroboration.
  \item[Low] sparse, indirect, or conflicting evidence. The tool flags these
  results explicitly in its output.
\end{description}

\textbf{Block confidence} is a property of an \emph{assessment}, expressing how
completely the user filled in the inputs that a given output block depends on.
It is computed as described below.

A typology may have high evidence confidence in an assessment with low block
confidence, or the reverse. Reporting a single merged figure would lose the
distinction between ``we know this intervention well but you told us little
about your site'' and ``you described your site thoroughly but the evidence for
this intervention is thin''. Those call for different responses.

\subsection{Branched confidence}
\label{sec:confidence}

A single confidence rating would be dominated by whichever input block happened
to be least complete, and would obscure the common case of a site with good
physical data and no economic data. Confidence is therefore computed separately
for each output block --- cooling, energy, economic, and equity --- from the
completeness of the inputs feeding that block:

\begin{table}[H]
\centering
\caption[The confidence model]{The block confidence model. Completeness is the
proportion of the fields relevant to that output block which the user
supplied.}
\label{tab:confidence}
\small
\begin{tabular}{@{}lll@{}}
\toprule
\textbf{Completeness of relevant fields} & \textbf{Confidence} &
\textbf{How the result should be read} \\
\midrule
Below \SI{40}{\percent}   & Low    & A screening signal for further investigation \\
\SIrange{40}{70}{\percent} & Medium & Indicative; the main drivers are known \\
Above \SI{70}{\percent}   & High   & As reliable as a screening instrument gets \\
\bottomrule
\end{tabular}
\end{table}

Two additional constraints apply.

\subsubsection{Evidence confidence propagates}
\label{sec:evidence-propagation}

A typology rated \emph{low} evidence confidence in the library --- green
fa\c{c}ade, rain garden / bioswale, and courtyard greening at version
\methver{} --- caps the cooling-block confidence at \emph{medium} regardless of
how complete the user's inputs are.

The reason is straightforward: complete inputs cannot compensate for thin
underlying evidence. A user who answers every question about a site where a
green fa\c{c}ade is proposed has told the tool a great deal about the site and
nothing at all about whether green fa\c{c}ades reliably reduce street-level air
temperature, which is the question the disagreement between
\textcite{keravec2026} and \textcite{zolch2016} leaves open. The interface
states the reason for the cap explicitly rather than showing an unexplained
ceiling.

\subsubsection{A high score with low confidence is a flag, not a verdict}

The combination of a high Opportunity Score and low block confidence is
presented as an invitation to investigate further, never as a conclusion. In
practice this is the most dangerous output the tool can produce, because the
headline number is exactly what gets extracted and quoted. The interface and
the exported report attach the confidence rating to the score wherever the
score appears, and \Cref{sec:misuse} identifies quoting a score without its
confidence as a misuse.

\subsection{Missing data}
\label{sec:missing-data}

Required inputs must be present; the assessment does not proceed without them.
Optional inputs, when absent, follow explicit rules:

\begin{itemize}
  \item \textbf{Qualitative inputs} default to \emph{unknown}, which normalises
  to the neutral 50 of \Cref{sec:unknown}.
  \item \textbf{Quantitative inputs} propagate as \emph{not calculated}. They
  are \textbf{never} treated as zero. A zero would understate the result while
  appearing to be a finding --- an energy saving reported as
  \qty{0}{\kilo\watt\hour} reads as ``this intervention saves no energy'',
  which is a claim, whereas the truth is ``we were not told the demand'', which
  is an absence.
  \item \textbf{Derived outputs} whose inputs are missing carry an explicit
  status naming the reason, as the energy statuses of \Cref{sec:energy}
  illustrate.
\end{itemize}

Every default the engine applies is itemised in the result and reproduced in
the exported report, so that a reader can see exactly which numbers came from
the user and which from the tool. This list is not a diagnostic aid; it is part
of the output, and a report circulated without it has removed the information
needed to judge the numbers it contains.
